\section{Introduction}
\label{sec:intro}

\IEEEPARstart{M}{ultimodal} classification with hyperspectral imagery (HSI) and LiDAR data plays an important role in remote sensing applications such as urban planning, forest monitoring, and geological mapping.
HSI provides dense spectral measurements for material discrimination, whereas LiDAR contributes height and structural cues that help distinguish spectrally similar land-cover classes, making the two modalities complementary~\cite{hong2022multimodal}.
Building on this complementarity, recent HSI-LiDAR fusion networks have achieved strong performance in multimodal remote sensing classification with feature-level and decision-level fusion strategies~\cite{hang2020coupled,fang2022s2enet,li2022mahidfnet}.

With the rapid growth of remote sensing data and the increasing diversity of sensing platforms, land-cover observations accumulate over time across new regions, sensors, and acquisition conditions, giving rise to a lifelong learning scenario in which both the label space and the data domain evolve after initial deployment.
Continual learning studies how models can learn from sequentially arriving data while retaining knowledge acquired from previous tasks~\cite{delange2021continualsurvey,wang2024continualsurvey}.
Typical continual learning settings~\cite{vandeven2022three} include class-incremental learning (CIL), where new categories arrive over time and the model predicts over all classes seen so far, typically within a shared data domain~\cite{rebuffi2017icarl,hou2019lucir}; domain-incremental learning (DIL), where the input distribution shifts across tasks while the label space and task structure remain unchanged~\cite{wang2024pina}; and task-incremental learning (TIL), where the stream is partitioned into tasks and the task identity is provided at inference~\cite{serra2018hat}.
This paper focuses on cross-domain class-incremental learning (CDCIL), where class-incremental label growth and domain-incremental distribution shift are coupled within a single task stream, as illustrated in Fig.~\ref{fig:cdcil_setting}.
  
\begin{figure*}[t!]
  \centering
  \includegraphics[width=\textwidth]{figures/fig1/fig1.pdf}
  \caption{\textbf{Cross-domain class-incremental learning setting on HSI and LiDAR data.}
  The model is trained sequentially on $9$ tasks across three sensors (MUUFL, Trento, and Houston 2013), with the class set accumulating across tasks ($4 \to 32$ classes total). Each task introduces a disjoint subset of land-cover classes; the cumulative class set after task $t$ is denoted ``Class set $t$''.}
  \label{fig:cdcil_setting}
\end{figure*}

CDCIL inherits the core difficulties of both CIL and DIL, but their coupling creates new challenges that neither setting addresses in isolation.
As shown in Fig.~\ref{fig:cdcil_setting}, each domain is divided into a sequence of class-incremental tasks, where every task introduces a subset of new land-cover classes and the cumulative label space grows within that domain.
After the tasks of one domain are completed, the stream moves to a different sensor and location; changes in wavelength range, band count, and ground sample distance (GSD) shift the input distribution while class accumulation continues.
The model therefore faces a sharpened stability--plasticity dilemma~\cite{mermillod2013stability}: it needs to avoid catastrophic forgetting of old classes from previously seen domains, absorb new classes, and align feature representations across changing domains simultaneously.
This forgetting risk arises at two levels: intra-domain forgetting, where decision boundaries fitted to new classes reduce the separability of earlier classes from the same domain; and inter-domain forgetting, where adaptation to a new sensor or location shifts these boundaries away from earlier-domain classes.
These boundary shifts are amplified by semantic overlap across domains.
For example, Fig.~\ref{fig:cdcil_setting} shows that the class ``Trees'' appears in both MUUFL and Houston 2013, ``Apple Trees'' appears in Trento, and ``Road'' appears in all three domains.
The feature space therefore has to support discrimination at two levels: the model has to separate domains with different sensing conditions while distinguishing semantically related classes within and across those domains.

Current methods addressing catastrophic forgetting can be broadly grouped into five directions: replay-based~\cite{rebuffi2017icarl,hou2019lucir,douillard2020podnet,wang2022foster}, regularization-based~\cite{kirkpatrick2017ewc,li2017lwf}, architecture-based~\cite{wang2022l2p,wang2022dualprompt,smith2023codaprompt,liang2024inflora}, meta-learning-based~\cite{yang2026pmpt,zhu2025pass}, and feature-space-based~\cite{petit2023fetril,mcdonnell2023ranpac,gomezvilla2024ldc}.
Remote-sensing CIL methods extend these strategies to hyperspectral and multimodal classification scenarios, but are mostly evaluated within a single scene, sensor, or dataset~\cite{deng2024dcprn,xu2025crossacl,zhang2024jfi_acl,li2026dfpem,yang2026pmpt}.
CDCIL has also attracted attention in other fields, such as medical diagnosis and industrial fault diagnosis, where new categories arrive under changing domains~\cite{yang2023cdfscil,shi2024cdcibn}.
In remote sensing, the closest existing settings sit one or two axes away from ours: HyperKD~\cite{li2025hyperkd} performs single-modality HSI lifelong learning with each dataset as one coarse task; CDCIRSC~\cite{zhang2024cdcirssc} targets scene-level cross-domain CIL on different data sources; and recent HSI-LiDAR transfer methods~\cite{liu2025prompt} adapt across datasets through joint rather than incremental training.
Multimodal HSI-LiDAR CDCIL with finer per-domain task granularity therefore exposes representation-level difficulties that remain insufficiently captured by existing settings, especially when class boundaries evolve within each domain while sensor and scene conditions shift across domains.
This leaves open how to preserve separability between old and new classes within each domain, prevent forgetting across domains, and maintain HSI-LiDAR representations that distinguish semantically similar classes while accounting for domain-specific sensing conditions.

To diagnose these representation-level difficulties, we examine branch-level drift in representative HSI-LiDAR fusion backbones, including FusAtNet~\cite{mohla2020fusatnet}, S2ENet~\cite{fang2022s2enet}, Coupled CNNs~\cite{hang2020coupled}, MAHiDFNet~\cite{li2022mahidfnet}, and HCT~\cite{ghosh2023hct}.
We find that, across these backbones, the spatial-to-spectral drift ratio ranges from $0.60\times$ for FusAtNet to $5.66\times$ for HCT.
Coupled fusion architectures such as FusAtNet show weak or no spectral--spatial drift separation, whereas backbones with more explicit spectral--spatial separation, such as HCT, tend to allocate more drift to the spatial branches.
Motivated by this observation, instead of suppressing branch drift uniformly, we design \backbonename{}, a tri-branch backbone that isolates a capacity-bounded spectral stream as a stable anchor across domains while letting the spatial branches absorb domain-specific variation.
Within this pool, \backbonename{} attains a spectral drift of $0.081$ and a $7.05\times$ spatial-to-spectral ratio, supporting the intended anchor behavior.
Building on this anchor, \methodname{} is a lightweight two-layer CDCIL framework that trains the backbone during first-domain warm-up and keeps it frozen; subsequent domains are processed by forward feature extraction, a compact frozen-feature memory, and incremental updates of domain-specific heads.
The spectral branch routes each sample to a domain-specific head, while the spatial branches contribute drift-suppressed features for fine-grained classification within the routed domain.

The main contributions of this paper are summarized as follows:
\begin{itemize}
    \item We propose \methodname{}, a lightweight two-layer framework specifically tailored to multimodal HSI-LiDAR cross-domain class-incremental learning. It trains the backbone during first-domain warm-up and keeps it frozen; subsequent domains are processed by forward feature extraction, a compact frozen-feature memory, and incremental updates of domain-specific heads. Stable spectral features distinguish domains, whereas drift-suppressed spatial features support fine-grained class discrimination, thereby mitigating intra-domain and inter-domain forgetting under a sharpened stability--plasticity dilemma.
    \item \methodname{} is built on \backbonename{}, a tri-branch backbone motivated by a branch-level drift analysis across five representative HSI-LiDAR fusion backbones. The analysis suggests two design choices associated with stronger spectral--spatial drift asymmetry---isolating the spectral stream from cross-modal coupling and bounding its capacity---and \backbonename{} instantiates both, providing the low-drift spectral anchor used by \methodname{} for domain routing.
    \item We conduct comprehensive CDCIL experiments on MUUFL, Trento, and Houston 2013 with $9$ tasks and $32$ classes, covering baseline comparisons, branch-level drift analysis, spectral-routing evaluation, memory-footprint analysis, and ablation studies to verify the roles of the spectral anchor and drift-suppressed spatial features.
\end{itemize}
